\documentclass[10pt]{article}
\usepackage[margin=1in]{geometry}
\usepackage{graphicx,booktabs,amsmath,amssymb,hyperref,xcolor,caption,microtype,multirow}
\hypersetup{colorlinks=true,linkcolor=blue!50!black,urlcolor=blue!50!black,citecolor=blue!50!black}
\captionsetup{font=small,labelfont=bf}
\newcommand{\Lstar}{L$^{*}$}
\newcommand{\ci}[2]{[#1,\,#2]}
\newcommand{\todo}[1]{\textcolor{red}{[#1]}}

\title{Drift to Default: Long Autoregressive Video Rollouts Relax\\ to a Prompt-Set State}
\author{Yilin Wang\\ University of Iowa}
\date{}

\begin{document}
\maketitle

\begin{abstract}
Autoregressive video diffusion models generate long videos chunk by chunk, conditioned on their own output, and drift
over tens of seconds. We show that this drift is not a random walk but \emph{relaxation}. The colour and texture
statistics of a rollout move from those of its opening toward a default set by the text prompt, and the long-run state
is an equilibrium between two pulls:
\[
x_\infty \approx \alpha\, x_{\text{open}} + (1-\alpha)\,\mu_{\text{prompt}}.
\]
\begin{itemize}
\item \textbf{Generality.} About 22--40\% of an opening's deviation survives 31\,s. This holds for brightness, colour,
saturation, contrast and sharpness, across \todo{9} model configurations spanning DMD, SiD, GAN, LongLive, CausVid,
Rolling Forcing and SkyReels diffusion forcing. The default $\mu_{\text{prompt}}$ moves with the lighting described in
the prompt (16--45 \Lstar) and differs between models for the same words.
\item \textbf{Mechanism.} An attention sink keeps the opening in the context and sets $\alpha$. Fitted to the observed
persistence curves, a linear chunk-recursion model recovers the sink's per-chunk weight that we measure directly by
bit-exact KV-cache replay (0.05). The sink's tiny direct weight still sets the equilibrium because the prompt's
per-chunk pull is equally small.
\item \textbf{Two levers.} Raising $\alpha$ with a training-free, content-aware correction reduces drift and keeps
83--93\% of legitimate brightness changes; post-hoc colour matching keeps 7--12\%. Moving $\mu_{\text{prompt}}$ toward
the opening with a lighting description calibrated to measured brightness reduces drift by 26\% on Self Forcing; a
free-form VLM caption does not.
\end{itemize}
All experiments use paired seeds on a pinned GPU architecture. Six analyses were preregistered with hashed code, and we
report the two that failed.
\end{abstract}

\section{Introduction}
Causal (autoregressive) video diffusion models~\cite{causvid,selfforcing} make streaming, open-ended generation
possible and underlie interactive world models. Their known weakness is drift: over tens of seconds the picture
darkens, saturates, loses texture or collapses. Many recent methods target this failure, through rolling windows,
attention sinks, re-aligned positions, longer self-rollouts and test-time
correction~\cite{longlive,rollingforcing,deepforcing,pathwise,freqforcing}. They are usually evaluated with one sample
per prompt, an aggregate quality score and no uncertainty. What drift actually \emph{is} has not been measured.

We treat drift as a measurement problem and borrow the design tools of clinical research: paired seeds, controlled
interventions, held-out prompts, cluster-bootstrap intervals and preregistered analyses. Our central finding is
simple. A long rollout \emph{relaxes}: every low-level statistic moves from its opening value toward a default set by
the prompt, and a sink holds back a fixed share of the opening. This one picture explains several observations that
otherwise look unrelated:
\begin{itemize}
\item why sinks reduce drift;
\item why real camera frames ``drift more'' than generated ones;
\item why a test-time correction that copies the anchor's tone works but over-corrects at early denoising steps;
\item why describing the opening in the prompt fails while calibrating the lighting words succeeds.
\end{itemize}

\paragraph{Contributions.}
\begin{enumerate}
\item A measurement protocol for long-video drift: self-referenced metrics, held and \emph{perturbed} openings, paired
seeds pinned to one GPU architecture, and preregistration. It also exposes three evaluation pitfalls:
  \begin{itemize}
  \item the GPU architecture changes bf16 outputs;
  \item resumed jobs silently break pairing;
  \item a standard image-quality metric is blind to drift.
  \end{itemize}
\item The relaxation law, $x_\infty \approx \alpha x_{\text{open}} + (1-\alpha)\mu_{\text{prompt}}$. We establish it
interventionally, by shifting the brightness and colour of the same opening, across \todo{9} model configurations and
six statistics.
\item A quantitative mechanism. Long-horizon persistence curves and direct KV-cache interventions give consistent
per-chunk weights, which explains how a sink with a tiny weight sets the equilibrium.
\item Two practical levers matching the two terms of the law, both evaluated in preregistered tests against
baselines: a content-aware correction raises $\alpha$, and calibrated prompts move $\mu$.
\end{enumerate}
No model is trained. Everything is inference on public weights, so the protocol applies to any open autoregressive
video model.

\section{Related work}
\paragraph{Autoregressive video diffusion.}
\begin{itemize}
\item Diffusion Forcing~\cite{diffusionforcing} trains with independent per-frame noise.
\item CausVid~\cite{causvid} distils a bidirectional teacher~\cite{wan} into a causal few-step student with
DMD~\cite{dmd}.
\item Self Forcing~\cite{selfforcing} trains on the model's own rollouts.
\item Long-video successors add rolling KV windows, attention sinks (as in streaming language
models~\cite{streamingllm}), RoPE re-alignment, LoRA-based extension, and joint rolling
denoising~\cite{longlive,rollingforcing,deepforcing}.
\item SkyReels-V2~\cite{skyreels} extends videos by diffusion forcing over overlapping segments.
\end{itemize}

\paragraph{Error accumulation and correction.} Meta-ARVDM~\cite{metaarvdm} analyses error accumulation theoretically.
Pathwise test-time correction~\cite{pathwise} and FreqForcing~\cite{freqforcing} correct drift at inference; our
correction is a deliberately minimal member of this family. It transfers only the DC tone at one early step, which
lets us study the step as a knob.

\paragraph{Evaluation.} VBench~\cite{vbench} scores imaging quality with MUSIQ~\cite{musiq}. We show that MUSIQ is
nearly blind to drift, and we measure drift relative to each prompt's own fresh-start distribution, in colour
statistics and in DINOv2~\cite{dinov2} features.

\section{Setup}\label{sec:setup}
\paragraph{Models.} All configurations generate 832$\times$480 video at 16\,fps in chunks of 3 latent frames (0.75\,s),
except SkyReels (960$\times$544, 24\,fps); details are in Table~\ref{tab:models}.
\begin{itemize}
\item \textbf{Self Forcing~\cite{selfforcing} checkpoints:} DMD (main), SiD, SiD-v2, GAN, and the 10-second variant.
All share one architecture (Wan2.1-1.3B~\cite{wan}) and differ only in the distillation objective.
\item \textbf{CausVid~\cite{causvid}:} autoregressive DMD checkpoint.
\item \textbf{LongLive~\cite{longlive}:} LoRA, official window 12 + sink 3.
\item \textbf{Rolling Forcing~\cite{rollingforcing}:} rolling joint denoising with a fixed anchor block.
\item \textbf{SkyReels-V2-DF-1.3B~\cite{skyreels}:} diffusion forcing over overlapping segments.
\end{itemize}

\paragraph{Held and perturbed openings.} To separate drift from what the prompt asks for, an opening image is held for
0.56\,s, encoded as the first chunk (or as the prefix video for SkyReels), and continued for 31\,s with a generic driving
prompt. Openings are NAVSIM~\cite{navsim} front-camera frames or first frames generated by the model. In the
interventional study, the same opening is shifted in CIELAB before encoding: \Lstar{} by $\pm$8 and $\pm$16, and
$b^*$ by $\pm$12.

\paragraph{Metrics.} For each video we record per-frame Lab means, saturation, \Lstar{} contrast and Laplacian
sharpness. The \emph{persistence} $\alpha$ is the within-scene slope of the end value (last 2\,s) on the opening value.
For prompt-only rollouts, drift is measured against 10 fresh 5-second generations of the same prompt: $D$ is the
Mahalanobis distance of colour statistics and $K$ the DINOv2 nearest-neighbour distance.

\paragraph{Statistics.} Seeds are paired across arms, and every arm of a scene runs in one job on one GPU
architecture. We use mixed models or seed-averaged paired differences, with scene-cluster bootstrap intervals and Holm
correction within each preregistered family. Each preregistration froze its code by SHA-256 hash before generation, and
the last four were committed publicly first.

\section{Drift is relaxation toward a prompt-set default}\label{sec:relax}
\paragraph{Observational evidence.} Across 40 held openings with a generic prompt, the change of every statistic over
31\,s is strongly anti-correlated with its opening value (Table~\ref{tab:attr}): $r$ between $-0.79$ and $-0.95$ for
\Lstar, $a^*$, $b^*$, saturation, sharpness and contrast, on both Self Forcing and LongLive. For brightness:
\[
\Delta L = 22.8 - 0.69\,L_{\text{open}}, \quad R^2 = 0.82.
\]
\begin{itemize}
\item Only 22--40\% of an opening's deviation survives.
\item Real and generated openings fall on the same line ($p=0.80$), so real frames ``drift more'' only because they are
brighter.
\end{itemize}

\paragraph{The default is set by the prompt.} Continuing the same openings under four lighting prompts gives four
levels (Fig.~\ref{fig:default}b). On Self Forcing these are sunny 43, overcast 34, dusk 19 and night 16 \Lstar; the
lighting term has $p\approx5\times10^{-78}$. On LongLive the same words give 45, 40, 32 and 21. The prompt sets the
target, and the mapping from words to brightness is model-specific.

\paragraph{Interventional evidence across models.} \todo{Relaxation study (preregistered): 24 scenes $\times$ 7
opening versions $\times$ 9 configurations. Within-scene $\alpha_L$, $\alpha_b$ per configuration; H1--H4; table of
$\alpha$, half-life and default level per model; fresh-start comparison; long-horizon subset.}

\section{Mechanism: an equilibrium between the anchor and the prompt}\label{sec:mech}
Model a chunk's mean statistic as a linear recursion:
\[
x_{k+1}-\mu = w_w\big(\bar x_{\text{window}}-\mu\big) + w_s\big(x_{\text{sink}}-\mu\big).
\]
The opening counts as a window chunk until it rolls out; after that, with a sink, it keeps a direct pull $w_s$. The
prompt's pull is $w_p = 1-w_w-w_s$. Without a sink the opening is eventually forgotten ($\alpha\to0$). With a sink the
equilibrium share is
\[
\alpha_\infty = \frac{w_s}{w_s + w_p}.
\]
The share is large whenever the prompt's pull is as weak as the sink's.

\paragraph{Two independent measurements agree.} Bit-exact KV-cache replay lets us shift only the sink or only the
window and regenerate the next chunk from identical noise. It measures $w_s = 0.053$ \ci{0.040}{0.066} for synthetic
tone shifts, $0.036$ in a dose-response design, and $w_w \approx 0.76$ (synthetic) to $0.90$--$0.99$ (natural
differences). Fitting the recursion to the long-horizon persistence curve alone gives:
\begin{itemize}
\item Self Forcing: $w_s = 0.049$ \ci{0.005}{0.069}, $w_w = 0.85$, $\alpha_\infty = 0.32$;
\item LongLive: $w_s = 0.054$, $w_w = 0.81$.
\end{itemize}
This resolves an earlier puzzle: the sink's direct effect on the next chunk is tiny, yet it roughly halves drift.
\todo{Interventional test with sink 0 vs 3 from the relaxation study.}

\section{Two levers}\label{sec:levers}
\paragraph{Raising $\alpha$: a content-aware correction.} At denoising step 1 of 4, the noisy chunk is run once more
with a temporary KV cache that holds only the anchor. The difference between the per-channel spatial means (the DC tone)
of that prediction and of the real one is added back, scaled by $\beta$. Nothing else changes, including the main cache
and the random stream.
\begin{itemize}
\item \textbf{Confirmation (preregistered, 40 held-out prompts $\times$ 4 seeds).} Colour drift $D$ falls by 1.48
\ci{1.19}{1.76} at $\beta=1$ and by 1.78 at $\beta=1.5$. Structure drift $K$ is non-inferior. On held-out prompts with
legitimate brightness changes, 83\% / 93\% of the change is kept.
\item \textbf{Baselines.} Post-hoc colour matching to the anchor (mean, or Reinhard mean and standard deviation)
reaches similar $D$ ($-0.89$, $-1.58$) but keeps only 7\% / 12\% of legitimate transitions (Table~\ref{tab:baselines}).
\item \textbf{The step is a knob.} Correcting at step 0 removes more drift but erases tunnels (56\% kept); steps 2--3
do nothing.
\item \textbf{Cost and transfer.} The correction costs one forward pass per chunk and some flicker. On LongLive it
reduces $D$ by half as much.
\end{itemize}

\paragraph{Moving $\mu$: calibrated prompts.}
\begin{itemize}
\item \textbf{Free-form caption (preregistered).} Describing the opening with a Qwen2.5-VL~\cite{qwenvl} caption does
not reduce brightness drift on either model (Self Forcing $-1.21$ \ci{-3.09}{0.36}). The caption describes half of the
normally lit frames as ``dim'' and pulls them darker. It does improve content retention (+0.10 / +0.15).
\item \textbf{Calibrated lighting (preregistered).} Choosing the lighting sentence whose measured default is nearest
the opening's brightness reduces $|\Delta L|$ by 26\% on Self Forcing (Holm $p = 0.013$); with the correction, by 44\%.
\item \textbf{Calibration is per model.} On LongLive, the preregistered test was not significant even with LongLive's
own calibration; the lighting-only prompt reduces drift by 34\% (secondary).
\end{itemize}
\todo{Out-of-sample prediction of these outcomes from $\alpha$ and $\mu$.}

\section{Context strategies}\label{sec:ctx}
A preregistered factorial study (10 strategies, 40 prompts $\times$ 4 seeds) found:
\begin{itemize}
\item a 3-frame sink cuts colour drift by 7.56 \ci{6.81}{8.32};
\item window size matters only without a sink;
\item RoPE re-alignment, which Self Forcing was not trained for, does not help;
\item MUSIQ varies only between 66 and 69 across strategies whose drift differs threefold;
\item the GPU architecture changes outputs, so 70\% of nominally paired comparisons were not paired.
\end{itemize}

\section{What did not hold}
We report claims that were revised or withdrawn after checks:
\begin{itemize}
\item that MUSIQ rewards collapsed videos (it is merely insensitive);
\item that the sink stabilises the default (trend difference $+0.5$ \ci{-2.4}{3.5});
\item that the model follows its context by 75\% (only for artificial shifts);
\item that early dynamics forecast the final brightness;
\item that describing the opening suffices.
\end{itemize}

\section{Limitations}
\begin{itemize}
\item All models but SkyReels share the Wan2.1-1.3B base.
\item Openings are held images, not real video.
\item The correction acts on colour only.
\item Lighting calibration uses four discrete levels and cannot match openings brighter than the sunny default.
\end{itemize}

\begin{thebibliography}{99}\small
\bibitem{causvid} T. Yin, Q. Zhang, R. Zhang, W. T. Freeman, F. Durand, E. Shechtman, X. Huang. From slow bidirectional to fast autoregressive video diffusion models. CVPR 2025 (arXiv:2412.07772).
\bibitem{selfforcing} X. Huang, Z. Li, G. He, M. Zhou, E. Shechtman. Self Forcing: Bridging the train-test gap in autoregressive video diffusion. arXiv:2506.08009, 2025.
\bibitem{diffusionforcing} B. Chen, D. Mart\'i Mons\'o, Y. Du, M. Simchowitz, R. Tedrake, V. Sitzmann. Diffusion Forcing: Next-token prediction meets full-sequence diffusion. NeurIPS 2024 (arXiv:2407.01392).
\bibitem{wan} Wan Team. Wan: Open and advanced large-scale video generative models. arXiv:2503.20314, 2025.
\bibitem{dmd} T. Yin et al. One-step diffusion with distribution matching distillation. CVPR 2024 (arXiv:2311.18828).
\bibitem{sid} M. Zhou et al. Score identity distillation. ICML 2024 (arXiv:2404.04057).
\bibitem{longlive} S. Yang et al. LongLive: Real-time interactive long video generation. 2025.
\bibitem{rollingforcing} K. Liu et al. Rolling Forcing: Autoregressive long video diffusion in real time. 2025.
\bibitem{deepforcing} Deep Forcing. arXiv:2512.05081, 2025.
\bibitem{skyreels} SkyReels Team. SkyReels-V2: Infinite-length film generative model. arXiv:2504.13074, 2025.
\bibitem{streamingllm} G. Xiao, Y. Tian, B. Chen, S. Han, M. Lewis. Efficient streaming language models with attention sinks. ICLR 2024 (arXiv:2309.17453).
\bibitem{pathwise} Pathwise test-time correction for autoregressive video diffusion. arXiv:2602.05871, 2026.
\bibitem{freqforcing} FreqForcing. arXiv:2607.27110, 2026.
\bibitem{metaarvdm} Meta-ARVDM. arXiv:2503.10704, 2025.
\bibitem{vbench} Z. Huang et al. VBench: Comprehensive benchmark suite for video generative models. CVPR 2024 (arXiv:2311.17982).
\bibitem{musiq} J. Ke, Q. Wang, Y. Wang, P. Milanfar, F. Yang. MUSIQ: Multi-scale image quality transformer. ICCV 2021.
\bibitem{dinov2} M. Oquab et al. DINOv2: Learning robust visual features without supervision. TMLR 2024 (arXiv:2304.07193).
\bibitem{qwenvl} Qwen Team. Qwen2.5-VL technical report. arXiv:2502.13923, 2025.
\bibitem{navsim} D. Dauner et al. NAVSIM: Data-driven non-reactive autonomous vehicle simulation and benchmarking. NeurIPS 2024.
\bibitem{reinhard} E. Reinhard, M. Ashikhmin, B. Gooch, P. Shirley. Color transfer between images. IEEE CG\&A, 2001.
\end{thebibliography}
\end{document}
